\documentclass[11pt]{article}
\usepackage[utf8]{inputenc}
\usepackage{amsmath, amssymb, amsfonts}
\usepackage{hyperref}
\usepackage{listings}
\usepackage{xcolor}

\definecolor{codegreen}{rgb}{0,0.6,0}
\definecolor{codegray}{rgb}{0.5,0.5,0.5}
\definecolor{codepurple}{rgb}{0.58,0,0.82}
\definecolor{backcolour}{rgb}{0.95,0.95,0.92}

\lstset{
    language=Python,
    backgroundcolor=\color{backcolour},   
    commentstyle=\color{codegreen},
    keywordstyle=\color{magenta},
    numberstyle=\tiny\color{codegray},
    stringstyle=\color{codepurple},
    basicstyle=\ttfamily\small,
    breakatwhitespace=false,         
    breaklines=true,                 
    captionpos=b,                    
    keepspaces=true,                 
    numbers=left,                    
    numbersep=5pt,                  
    showspaces=false,                
    showstringspaces=false,
    showtabs=false,                  
    tabsize=2
}

\title{Spectral Solver for the Poisson-Beltrami Problem on $\mathbb{T}^3$}
\author{Gemini CLI Solver Documentation}
\date{\today}

\begin{document}

\maketitle

\section{Problem Statement}
We consider the Poisson problem for the Laplace-Beltrami operator $\Delta_g$ on a triply periodic domain $\Omega = [0, 2\pi)^3$:
\begin{equation}
    \Delta_g u = f
\end{equation}
where $g_{ij}$ is a given metric tensor field. The Laplace-Beltrami operator is defined as:
\begin{equation}
    \Delta_g u = \frac{1}{\sqrt{g}} \partial_i \left( \sqrt{g} g^{ij} \partial_j u \right)
\end{equation}
where $g = \det(g_{ij})$ and $g^{ij}$ is the inverse metric tensor.

\section{Numerical Strategy}
\subsection{Discretization}
The problem is solved using a \textbf{pseudospectral method}:
\begin{itemize}
    \item \textbf{Derivatives:} All spatial derivatives ($\partial_j u$ and the outer divergence $\partial_i$) are computed in Fourier space using the Fast Fourier Transform (FFT).
    \item \textbf{Variable Coefficients:} The product $V^i = \sqrt{g} g^{ij} (\partial_j u)$ is computed in physical space to handle the arbitrary spatial dependence of the metric.
\end{itemize}

\subsection{Linear Solver and Preconditioning}
For high resolutions (up to $1024^3$), the resulting linear system is large and potentially ill-conditioned. We provide two strategies:
\begin{itemize}
    \item \textbf{Spectral Laplacian Preconditioner:} A constant-coefficient approximation that is diagonal in Fourier space.
    \item \textbf{Multigrid V-Cycle:} A full hierarchical solver in \texttt{poisson\_beltrami\_multigrid.py} that uses spectral restriction and prolongation across grid levels.
\end{itemize}

\subsection{GPU Readiness}
The solver is implemented entirely using vectorized \texttt{numpy} operations (specifically \texttt{einsum} for tensor contractions). This structure allows for a near-seamless transition to \texttt{CuPy} for execution on NVIDIA GPUs by simply replacing the \texttt{numpy} and \texttt{scipy.fft} imports.

\section{Documentation}
\subsection{Classes}
\begin{itemize}
    \item \texttt{BeltramiPoissonSolver}: Standard solver with spectral Laplacian preconditioning (in \texttt{poisson\_beltrami\_3d.py}).
    \item \texttt{MultigridBeltramiSolver}: Advanced solver with full V-cycle multigrid (in \texttt{poisson\_beltrami\_multigrid.py}).
\end{itemize}

\subsection{Usage Example (Multigrid)}
\begin{lstlisting}[language=Python]
from poisson_beltrami_multigrid import MultigridBeltramiSolver
import numpy as np

# Grid size
N = 1024
solver = MultigridBeltramiSolver(N, N, N, n_levels=4)

# Define metric g_ij (3,3,N,N,N) array
# solver.set_metric(g_ij)
# u = solver.solve(f)
\end{lstlisting}

\end{document}
